% GENERATED FILE. Edit canonical lessons, then run npm run generate:books.
% canonical-source-hash: fnv1a64:bc5e848691bbf990
% canonical-lessons: TA-C549-mantapam, TA-C549-nikalcciniral, TA-C549-munnottam, TA-C549-ayattam, TA-C549-naluvu

\chapter{Planning a Function}
\label{ch:ta-planning-a-function}
\hlchaptermodality{549}

\begin{chapteropening}
\textbf{By the end of this chapter:} \emph{I can talk about booking a hall and planning an event.}

Complete the last lesson of chapter 549: I can talk about booking a hall and planning an event.
\end{chapteropening}

\section[\texorpdfstring{\ta{மண்டபம்} (maṇṭapam)}{மண்டபம் (maṇṭapam)}]{\ta{மண்டபம்} (maṇṭapam) — a function hall}

\label{lesson:TA-C549-mantapam}



\begin{quote}
Before the new one: say the Tamil for a fraud, then the Tamil for to break into.
\end{quote}

\subsection*{You'll want to know: \ta{மண்டபம்}}
\textbf{\ta{மண்டபம்}} — \emph{maṇṭapam} — \textquotedblleft{}a function hall\textquotedblright{}.

\begin{quote}
\ta{மண்டபம்} \ta{கிடைத்தால்} \ta{திருமணம்} \ta{அடுத்த} \ta{மாதம்} \ta{நடக்கும்}.
\end{quote}

\begin{quote}
\emph{maṇṭapam kiṭaittāl tirumaṇam aṭutta mātam naṭakkum.}
\end{quote}

\begin{quote}
If we get the hall, the wedding will be next month.
\end{quote}

Weddings in South India are often held in a \ta{திருமண} \ta{மண்டபம்}.

\begin{cousinweb}
\textbf{\ta{மண்டபம்}} (\emph{maṇṭapam}) is borrowed from Sanskrit \emph{maṇḍapa}, \textquotedblleft{}an open temporary shed, a pavilion\textquotedblright{} (CDIAL 9740).
\end{cousinweb}

\begin{grammarlens}[title={What you've built}]
One more word for saying what happens if.
\end{grammarlens}

\subsection*{Your turn}
\begin{itemize}
\raggedright
  \item \emph{Say it:} \emph{maṇṭapam}
  \item \emph{Say it:} \emph{maṇṭapam}, once more
  \item \emph{Say it:} \emph{maṇṭapam kiṭaittāl tirumaṇam aṭutta mātam naṭakkum.} — If we get the hall, the wedding will be next month.
  \item \emph{Recall:} say the Tamil for a fraud, then the Tamil for to break into, then say \emph{maṇṭapam} again
\end{itemize}

\begin{tcolorbox}[breakable,colback=teal!4,colframe=teal!35!black,title={Before you move on}]
What does \ta{மண்டபம்} mean? (\textquotedblleft{}A function hall\textquotedblright{}.) Say it once more.
\end{tcolorbox}
\section[\texorpdfstring{\ta{நிகழ்ச்சிநிரல்} (nikaḻcciniral)}{நிகழ்ச்சிநிரல் (nikaḻcciniral)}]{\ta{நிகழ்ச்சிநிரல்} (nikaḻcciniral) — an agenda; a programme of events}

\label{lesson:TA-C549-nikalcciniral}



\begin{quote}
Before the new one: say the Tamil for to break into, then the Tamil for a function hall.
\end{quote}

\subsection*{You'll want to know: \ta{நிகழ்ச்சிநிரல்}}
\textbf{\ta{நிகழ்ச்சிநிரல்}} — \emph{nikaḻcciniral} — \textquotedblleft{}an agenda; a programme of events\textquotedblright{}.

\begin{quote}
\ta{கூட்டத்தின்} \ta{நிகழ்ச்சிநிரலை} \ta{எல்லோருக்கும்} \ta{அனுப்பினேன்}.
\end{quote}

\begin{quote}
\emph{kūṭṭattiṉ nikaḻcciniralai ellōrukkum aṉuppiṉēṉ.}
\end{quote}

\begin{quote}
I sent the agenda of the meeting to everyone.
\end{quote}

Word for word \textquotedblleft{}a row of events\textquotedblright{}: \ta{நிகழ்ச்சி} (\emph{nikaḻcci}) + \ta{நிரல்} \textquotedblleft{}a row\textquotedblright{}.

\begin{cousinweb}
\textbf{\ta{நிகழ்ச்சிநிரல்}} (\emph{nikaḻcciniral}) is a native compound: \ta{நிகழ்ச்சி} (\emph{nikaḻcci}), \textquotedblleft{}an event\textquotedblright{}, and \ta{நிரல்} (\emph{niral}), \textquotedblleft{}an order, a row\textquotedblright{} (DEDR 3673).
\end{cousinweb}

\begin{grammarlens}[title={What you've built}]
One more word for saying what happens if.
\end{grammarlens}

\subsection*{Your turn}
\begin{itemize}
\raggedright
  \item \emph{Say it:} \emph{nikaḻcciniral}
  \item \emph{Say it:} \emph{nikaḻcciniral}, once more
  \item \emph{Say it:} \emph{kūṭṭattiṉ nikaḻcciniralai ellōrukkum aṉuppiṉēṉ.} — I sent the agenda of the meeting to everyone.
  \item \emph{Recall:} say the Tamil for to break into, then the Tamil for a function hall, then say \emph{nikaḻcciniral} again
\end{itemize}

\begin{tcolorbox}[breakable,colback=teal!4,colframe=teal!35!black,title={Before you move on}]
What does \ta{நிகழ்ச்சிநிரல்} mean? (\textquotedblleft{}An agenda; a programme of events\textquotedblright{}.) Say it once more.
\end{tcolorbox}
\section[\texorpdfstring{\ta{முன்னோட்டம்} (muṉṉōṭṭam)}{முன்னோட்டம் (muṉṉōṭṭam)}]{\ta{முன்னோட்டம்} (muṉṉōṭṭam) — a preview; a trial run}

\label{lesson:TA-C549-munnottam}



\begin{quote}
Before the new one: say the Tamil for a function hall, then the Tamil for an agenda.
\end{quote}

\subsection*{You'll want to know: \ta{முன்னோட்டம்}}
\textbf{\ta{முன்னோட்டம்}} — \emph{muṉṉōṭṭam} — \textquotedblleft{}a preview; a trial run\textquotedblright{}.

\begin{quote}
\ta{புதிய} \ta{படத்தின்} \ta{முன்னோட்டம்} \ta{நேற்று} \ta{வெளியானது}.
\end{quote}

\begin{quote}
\emph{putiya paṭattiṉ muṉṉōṭṭam nēṟṟu veḷiyāṉatu.}
\end{quote}

\begin{quote}
The trailer of the new film came out yesterday.
\end{quote}

Word for word \textquotedblleft{}a running ahead\textquotedblright{}: \ta{முன்} + \ta{ஓட்டம்}.

\begin{cousinweb}
\textbf{\ta{முன்னோட்டம்}} (\emph{muṉṉōṭṭam}) is a native compound: \ta{முன்} (\emph{muṉ}), \textquotedblleft{}in front, before\textquotedblright{} (DEDR 5020), and \ta{ஓட்டம்} (\emph{ōṭṭam}), \textquotedblleft{}a run\textquotedblright{}, from \ta{ஓடு}, \textquotedblleft{}to run\textquotedblright{} (DEDR 1041).
\end{cousinweb}

\begin{grammarlens}[title={What you've built}]
One more word for saying what happens if.
\end{grammarlens}

\subsection*{Your turn}
\begin{itemize}
\raggedright
  \item \emph{Say it:} \emph{muṉṉōṭṭam}
  \item \emph{Say it:} \emph{muṉṉōṭṭam}, once more
  \item \emph{Say it:} \emph{putiya paṭattiṉ muṉṉōṭṭam nēṟṟu veḷiyāṉatu.} — The trailer of the new film came out yesterday.
  \item \emph{Recall:} say the Tamil for a function hall, then the Tamil for an agenda, then say \emph{muṉṉōṭṭam} again
\end{itemize}

\begin{tcolorbox}[breakable,colback=teal!4,colframe=teal!35!black,title={Before you move on}]
What does \ta{முன்னோட்டம்} mean? (\textquotedblleft{}A preview; a trial run\textquotedblright{}.) Say it once more.
\end{tcolorbox}
\section[\texorpdfstring{\ta{ஆயத்தம்} (āyattam)}{ஆயத்தம் (āyattam)}]{\ta{ஆயத்தம்} (āyattam) — readiness; preparation}

\label{lesson:TA-C549-ayattam}



\begin{quote}
Before the new one: say the Tamil for an agenda, then the Tamil for a preview.
\end{quote}

\subsection*{You'll want to know: \ta{ஆயத்தம்}}
\textbf{\ta{ஆயத்தம்}} — \emph{āyattam} — \textquotedblleft{}readiness; preparation\textquotedblright{}.

\begin{quote}
\ta{விருந்தினர்} \ta{வந்தால்} \ta{எல்லாம்} \ta{ஆயத்தமாக} \ta{இருக்க} \ta{வேண்டும்}.
\end{quote}

\begin{quote}
\emph{viruntiṉar vantāl ellām āyattamāka irukka vēṇṭum.}
\end{quote}

\begin{quote}
If the guests come, everything must be ready.
\end{quote}

With -\ta{ஆக}: \ta{ஆயத்தமாக} \ta{இரு}, to be ready.

\begin{cousinweb}
The origin of \textbf{\ta{ஆயத்தம்}} (\emph{āyattam}) is not settled; one suggestion is Sanskrit \emph{āyatta}, \textquotedblleft{}being at someone's disposal\textquotedblright{} (Monier-Williams), with the Tamil ending -\ta{அம்}.
\end{cousinweb}

\begin{grammarlens}[title={What you've built}]
One more word for saying what happens if.
\end{grammarlens}

\subsection*{Your turn}
\begin{itemize}
\raggedright
  \item \emph{Say it:} \emph{āyattam}
  \item \emph{Say it:} \emph{āyattam}, once more
  \item \emph{Say it:} \emph{viruntiṉar vantāl ellām āyattamāka irukka vēṇṭum.} — If the guests come, everything must be ready.
  \item \emph{Recall:} say the Tamil for an agenda, then the Tamil for a preview, then say \emph{āyattam} again
\end{itemize}

\begin{tcolorbox}[breakable,colback=teal!4,colframe=teal!35!black,title={Before you move on}]
What does \ta{ஆயத்தம்} mean? (\textquotedblleft{}Readiness; preparation\textquotedblright{}.) Say it once more.
\end{tcolorbox}
\section[\texorpdfstring{\ta{நழுவு} (naḻuvu)}{நழுவு (naḻuvu)}]{\ta{நழுவு} (naḻuvu) — to slip away; to evade}

\label{lesson:TA-C549-naluvu}



\begin{quote}
Before the new one: say the Tamil for a preview, then the Tamil for readiness.
\end{quote}

\subsection*{You'll want to know: \ta{நழுவு}}
\textbf{\ta{நழுவு}} — \emph{naḻuvu} — \textquotedblleft{}to slip away; to evade\textquotedblright{}.

\begin{quote}
\ta{வேலை} \ta{வந்தபோது} \ta{அவன்} \ta{மெதுவாக} \ta{நழுவினான்}.
\end{quote}

\begin{quote}
\emph{vēlai vantapōtu avaṉ metuvāka naḻuviṉāṉ.}
\end{quote}

\begin{quote}
When work came along, he quietly slipped away.
\end{quote}

Past (I …): \ta{நழுவினேன்}, \emph{naḻuviṉēṉ}.

Also of a chance: \ta{வாய்ப்பு} \ta{நழுவியது}, the chance slipped away.

\begin{cousinweb}
\textbf{\ta{நழுவு}} (\emph{naḻuvu}) is native Dravidian: DEDR 3620 gives \ta{நழுவு} (\emph{naḻuvu}), \textquotedblleft{}to skulk away, to evade, to slip\textquotedblright{}.
\end{cousinweb}

\begin{grammarlens}[title={What you've built}]
One more word for saying what happens if.
\end{grammarlens}

\subsection*{Your turn}
\begin{itemize}
\raggedright
  \item \emph{Say it:} \emph{naḻuvu}
  \item \emph{Say it:} \emph{naḻuvu}, once more
  \item \emph{Say it:} \emph{vēlai vantapōtu avaṉ metuvāka naḻuviṉāṉ.} — When work came along, he quietly slipped away.
  \item \emph{Recall:} say the Tamil for a preview, then the Tamil for readiness, then say \emph{naḻuvu} again
\end{itemize}

\begin{tcolorbox}[breakable,colback=teal!4,colframe=teal!35!black,title={Before you move on}]
What does \ta{நழுவு} mean? (\textquotedblleft{}To slip away; to evade\textquotedblright{}.) Say it once more.
\end{tcolorbox}
